\subsection{Separable polynomials}

PROPOSITION 3. --- Let \(f\) be a non-zero polynomial in \(\mathbf{K}[\mathbf{X}]\) and let \(\Omega\) be an algebraically closed extension of \(\mathbf{K}\) . The following conditions are equivalent:

\begin{enumerate}
\def\labelenumi{\alph{enumi})}
\item
  The polynomial \(f\) is relatively prime to its derivative \(f'\) in \(\mathbf{K}[\mathbf{X}]\) .
\item
  Either \(\deg(f) = 0\) , or \(\deg(f) > 0\) and \(\operatorname{dis}(f) \neq 0\) (IV, p.~83).
\item
  There exists an extension \(\mathbf{L}\) of \(\mathbf{K}\) such that \(f\) splits in \(\mathbf{L}[\mathbf{X}]\) into a product of distinct polynomials of degree \(\leqslant 1\) .
\item
  The roots of \(f\) in \(\Omega\) are simple.
\item
  The K-algebra \(\mathbf{K}[\mathbf{X}] / (f)\) is etale (V, p.~28, Def. 1).
\end{enumerate}

\(a) \Rightarrow d)\) : Under the hypothesis \(a)\) there exist two polynomials \(g\) and \(h\) in K{[}X{]} such that \(fg + f'h = 1\) (IV, p.~12). Let \(a\) be a root of \(f\) in \(\Omega\) ; we have

\[
f ^ {\prime} (a) h (a) = f (a) g (a) + f ^ {\prime} (a) h (a) = 1,
\]

whence \(f'(a) \neq 0\) ; therefore \(a\) is a simple root of \(f\) in \(\Omega\) (IV, p.~17, Prop. 7). \(d) \Rightarrow c)\) : If \(d)\) holds, \(f\) splits in \(\Omega[X]\) into a product of distinct factors of degree \(\leqslant 1\) .

\(c) \Rightarrow b)\) : Under the hypothesis \(c)\) there exist an element \(\lambda \neq 0\) in L and distinct elements \(\alpha_{1}, \ldots, \alpha_{n}\) of L such that \(f(X) = \lambda(X - \alpha_{1}) \ldots (X - \alpha_{n})\) . If \(\deg(f) > 0\) , we have (IV, p.~83, Prop. 11),

\[
\operatorname{dis} (f) = \lambda^ {2 n - 2} \prod_ {i <   j} (\alpha_ {i} - \alpha_ {j}) ^ {2} \neq 0.
\]

\(b) \Rightarrow a)\) : Let \(c\) be the leading coefficient and D the discriminant of \(f\) ; the resultant of \(f'\) and \(f\) is equal to \(\pm cD\) (IV, p.~84, Formula (54)), hence is non-zero; therefore (IV, p.~78, Cor. 2) the polynomials \(f\) and \(f'\) are relatively prime in K{[}X{]}.

\(a) \Rightarrow e)\) : Let A be the K-algebra \(\mathbf{K}[\mathbf{X}] / (f)\) and \(x\) the image of X in A; by III, p.~573, Prop. 22, the A-module \(\Omega_{\mathrm{K}}(\mathbf{A})\) is generated by the elements \(dx\) , subject to the single relation \(f'(x) dx = 0\) . By V, p.~33, Th. 3, the K-algebra A is thus etale if and only if \(f'(x)\) is an invertible element in A, which means that \(f\) and \(f'\) are relatively prime in \(\mathbf{K}[\mathbf{X}]\) .

DEFINITION 2. --- A polynomial \(f \in \mathbf{K}[\mathbf{X}]\) is said to be separable if it is non-zero and satisfies the equivalent conditions a), b), c), d) and e) of Prop. 3.

Remarks. --- 1) Let L be an extension of K and f a non-constant polynomial in K{[}X{]}. By e) of Prop. 3 and V, p.~32, Cor. 2 it comes to the same to suppose that f is separable, whether considered as element of K{[}X{]} or of L{[}X{]}. On the other hand, it may well happen that f is irreducible in K{[}X{]} but not in L{[}X{]}.

\begin{enumerate}
\def\labelenumi{\arabic{enumi})}
\setcounter{enumi}{1}
\tightlist
\item
  Let \(f \in \mathbf{K}[\mathbf{X}]\) ; we know (IV, p.~13, Prop. 13) that there exist irreducible polynomials \(f_1, \ldots, f_m\) in \(\mathbf{K}[\mathbf{X}]\) such that \(f = f_1 \ldots f_m\) . Let \(\Omega\) be an algebraic closure of \(\mathbf{K}\) ; since an irreducible polynomial \(g \in \mathbf{K}[\mathbf{X}]\) is the minimal polynomial over \(\mathbf{K}\) of each of its roots in \(\Omega\) , two distinct irreducible polynomials in \(\mathbf{K}[\mathbf{X}]\) have no common root in \(\Omega\) . Condition \(d\) ) of Prop. 3 now shows that \(f\) is separable if and only if the polynomials \(f_1, \ldots, f_m\) are separable and pairwise distinct.
\end{enumerate}

PROPOSITION 4. --- Let \(f\) be an irreducible polynomial in \(\mathbf{K}[\mathbf{X}]\) . Then the following conditions are equivalent:

\begin{enumerate}
\def\labelenumi{\alph{enumi})}
\item
  \(f\) is separable.
\item
  There exists an extension \(\mathbf{L}\) of \(\mathbf{K}\) in which \(f\) has a simple root.
\item
  The derivative \(f'\) of \(f\) is not zero.
\item
  The field \(\mathbf{K}\) is of characteristic 0, or it is of characteristic \(p \neq 0\) and \(f \notin \mathbf{K}[\mathbf{X}^p]\) .
\end{enumerate}

We note first that an irreducible polynomial in K{[}X{]} is not constant. It is clear that a) implies b) (take an algebraic closure of K for L). If x is a simple root of f in an extension L of K, we have \(f'(x) \neq 0\) (IV, p.~17, Prop. 7), so b) implies c), and the equivalence of c) and d) follows from V, p.~9, Cor.

Suppose finally that \(f' \neq 0\) ; let \(x\) be a root of \(f\) in an algebraically closed extension \(\Omega\) of \(K\) . Since \(f\) is the minimal polynomial of \(x\) over \(K\) and \(\deg f' < \deg f\) , we have \(f'(x) \neq 0\) , and so \(x\) is a simple root of \(f\) (IV, p.~17, Prop. 7). Therefore \(f\) is separable and we have shown that \(c)\) implies \(a)\) .

COROLLARY 1. --- For a field K to be perfect it is necessary and sufficient that every irreducible polynomial of K{[}X{]} should be separable.

If the field K is of characteristic 0, K is perfect and every irreducible polynomial of K{[}X{]} is separable, by d) above. Suppose then that K has characteristic \(p \neq 0\) .

Suppose first that K is perfect. We have \(K[X^{p}] = K[X]^{p}\) , hence there exists no irreducible polynomial of K{[}X{]} belonging to \(K[X^{p}]\) . By Prop. 4, every irreducible polynomial of K{[}X{]} is then separable.

Suppose next that K is imperfect, whence \(K \neq K^{p}\) . Let a be an element of K not belonging to \(K^{p}\) ; the polynomial \(X^{p} - a\) is irreducible in K{[}X{]} (V, p.~24, Lemma 1), and it belongs to \(K[X^{p}]\) , and so is not separable.

COROLLARY 2. --- Let \(f \in \mathbf{K}[\mathbf{X}]\) be a non-zero polynomial. For \(f\) to be separable it is necessary and sufficient that there should exist an extension \(L\) of \(K\) which is a perfect field and such that \(f\) has no repeated factor in \(L[X]\) .

Let \(\Omega\) be an algebraic closure of \(\mathbf{K}\) ; if \(f\) is separable, \(f\) has no repeated factors in \(\Omega[\mathbf{X}]\) (Prop. 3, \(d\) )). Conversely, if \(\mathbf{L}\) is a perfect extension of \(\mathbf{K}\) such that \(f\) has no repeated factor in \(\mathbf{L}[\mathbf{X}]\) , then \(f\) is separable in \(\mathbf{L}[\mathbf{X}]\) (Cor. 1 and Remark 2), hence in \(\mathbf{K}[\mathbf{X}]\) (Remark 1).

